\documentclass[addpoints,spanish, 12pt,a4paper]{exam}
%\documentclass[answers, spanish, 12pt,a4paper]{exam}
 %\printanswers
\pointpoints{punto}{puntos}
\hpword{Puntos:}
\vpword{Puntos:}
\htword{Total}
\vtword{Total}
\hsword{Resultado:}
\hqword{Ejercicio:}
\vqword{Ejercicio:}

\usepackage[utf8]{inputenc}
\usepackage[spanish]{babel}
\usepackage{eurosym}
%\usepackage[spanish,es-lcroman, es-tabla, es-noshorthands]{babel}


\usepackage[margin=1in]{geometry}
\usepackage{amsmath,amssymb}
\usepackage{multicol}
\usepackage{yhmath}

\pointsinrightmargin % Para poner las puntuaciones a la derecha. Se puede cambiar. Si se comenta, sale a la izquierda.
\extrawidth{-2.4cm} %Un poquito más de margen por si ponemos textos largos.
\marginpointname{ \emph{\points}}

\usepackage{graphicx}

\graphicspath{{../img/}} 

\newcommand{\class}{4º Académicas}
\newcommand{\examdate}{\today}
\newcommand{\examnum}{Examen parcial 3ª evaluación}
\newcommand{\tipo}{A}


\newcommand{\timelimit}{50 minutos}

\renewcommand{\solutiontitle}{\noindent\textbf{Solución:}\enspace}


\pagestyle{head}
\firstpageheader{\includegraphics[width=0.2\columnwidth]{header_left}}{\textbf{Departamento de Matemáticas\linebreak \class}\linebreak \examnum}{\includegraphics[width=0.1\columnwidth]{header_right}}
\runningheader{\class}{\examnum}{Página \thepage\ of \numpages}
\runningheadrule


\begin{document}

\noindent
\begin{tabular*}{\textwidth}{l @{\extracolsep{\fill}} r @{\extracolsep{6pt}} }
\textbf{Nombre:} \makebox[3.5in]{\hrulefill} & \textbf{Fecha:}\makebox[1in]{\hrulefill} \\
 & \\
\textbf{Tiempo: \timelimit} & Tipo: \tipo 
\end{tabular*}
\rule[2ex]{\textwidth}{2pt}
Esta prueba tiene \numquestions\ ejercicios. La puntuación máxima es de \numpoints. 
La nota final de la prueba será la parte proporcional de la puntuación obtenida sobre la puntuación máxima. 

\begin{center}


\addpoints
 %\gradetable[h][questions]
	\pointtable[h][questions]
\end{center}

\noindent
\rule[2ex]{\textwidth}{2pt}

\begin{questions}


% \question[1] 
% \begin{solution} \end{solution}
% \addpoints

%\question 
%\begin{parts}
%\part[1]
%\begin{solution} \end{solution}
%\end{parts}
%\addpoints


%\question 
%\begin{parts}
%\part[1]
%\begin{solution} \end{solution}
%\end{parts}
%\addpoints

\question Calcular, usando las identidades fundamentales de la trigonometría, las razones trigonométricas de un ángulo agudo $\alpha$ sabiendo que:\begin{parts} \part[1] $\sen{\alpha}=\dfrac{\sqrt{2}}{2}$\begin{solution}  \\ $\sen{\alpha}=\dfrac{\sqrt{2}}{2}, \cos{\alpha}=\dfrac{\sqrt{2}}{2}, \tan{\alpha}=1$. \\ El ángulo agudo con esas razones es: $45^{\circ}$.\end{solution} 
\part[1] $\cos{\alpha}=\frac{\sqrt{1}}{2}$
\begin{solution} $\sin{x}=\frac{\sqrt{3}}{2}, \cos{x}=\frac{1}{2}, \tan{x}=\sqrt{3}$. \\ El ángulo agudo con esas razones es: $60^{\circ}$ \end{solution}
\part[1] $\tan{\alpha}=\sqrt{3}$\begin{solution}  \\ $\sen{\alpha}=\dfrac{\sqrt{3}}{2}, \cos{\alpha}=\dfrac{1}{2}, \tan{\alpha}=\sqrt{3}$. \\ El ángulo agudo con esas razones es: $60^{\circ}$.\end{solution} \end{parts} 

\question Calcular las razones trigonométricas de un ángulo $\alpha$ si:\begin{parts} \part[2] $\tan{\alpha}=\sqrt{3} \ y \cos{\alpha}<0$ \begin{solution}  \\ $\sin{\alpha}=- \dfrac{\sqrt{3}}{2}, \cos{\alpha}=- \dfrac{1}{2}, \tan{\alpha}=\sqrt{3}$. \\ El ángulo que cumple las condiciones del ejercicio es: $240^{\circ}$\end{solution} \part[2] $\tan{\alpha}=\dfrac{\sqrt{3}}{3} \ y \cos{\alpha}<0$ \begin{solution}  \\ $\sin{\alpha}=- \dfrac{1}{2}, \cos{\alpha}=- \dfrac{\sqrt{3}}{2}, \tan{\alpha}=\dfrac{\sqrt{3}}{3}$. \\ El ángulo que cumple las condiciones del ejercicio es: $210^{\circ}$\end{solution} \end{parts} 


\question Resuelve los triángulos rectángulos:\begin{parts} \part[1] Sabiendo que la hipotenusa mide 8 y un cateto 17 cm.\begin{solution} Los lados del triángulo miden: $8$, $15$, $17$ cm. Y los ángulos: $90.0 - 71.14 i$, $71.14 i$, $90$ º\end{solution} \part[1] Sabiendo que un cateto mide 15 cm. y su ángulo opuesto 30º\begin{solution} Los lados del triángulo miden: $15$, $25.98$, $30$ cm. Y los ángulos: $30$, $60$, $90$ º\end{solution} \part[1] Sabiendo que un cateto mide 30 cm. y su ángulo opuesto 45º\begin{solution} Los lados del triángulo miden: $30$, $30$, $42.43$ cm. Y los ángulos: $45$, $45$, $90$ º\end{solution} \part[1] Sabiendo que la hipotenusa mide 20 cm. y un ángulo 60º\begin{solution} Los lados del triángulo miden: $17.32$, $10$, $20$ cm. Y los ángulos: $60$, $30$, $90$ º\end{solution} \part[1] Sabiendo que un cateto mide 20 cm. y el ángulo opuesto al otro cateto 30º\begin{solution} Los lados del triángulo miden: $20$, $11.55$, $23.09$ cm. Y los ángulos: $60$, $30$, $90$ º\end{solution} \end{parts} 



\question Resuelve\begin{parts} \part[1] La diagonal menor de un rombo mide 20 cm y el ángulo menor es de 60$^{\circ}$. ¿Cuánto mide 
    la diagonal?¿Y el lado del rombo?\begin{solution} Miden $34.64$ y $20.0$ respectivamente\end{solution} \part[1] La diagonal menor de un rombo mide 40 cm y el ángulo menor es de 60$^{\circ}$. ¿Cuánto mide 
    la diagonal?¿Y el lado del rombo?\begin{solution} Miden $69.28$ y $40.0$ respectivamente\end{solution} \end{parts} 
    
    
\question Resuelve\begin{parts} \part[1] Un carpintero quiere construir una escalera de tijera cuyos brazos, una vez abiertos,
formen un ángulo de 60º. Si la altura de la escalera, estando abierta es de 2m, ¿qué
longitud deberá tener cada brazo?\begin{solution} $\dfrac{4 \sqrt{3}}{3} \to 2.31\ m $\end{solution} \part[1] Un carpintero quiere construir una escalera de tijera cuyos brazos, una vez abiertos,
formen un ángulo de 60º. Si la altura de la escalera, estando abierta es de 3m, ¿qué
longitud deberá tener cada brazo?\begin{solution} $2 \sqrt{3} \to 3.46\ m $\end{solution} \end{parts}
    
    
\question Resuelve\begin{parts} \part[1] Desde el punto donde estoy, la visual al punto más alto de una torre de 100 m que tengo 
        enfrente forma un ángulo de $30^{\circ}$ con la horizontal. ¿Cuántos m me tengo que acercar para que el ángulo
        sea de $60^{\circ}$?¿A cuántos metros estaba inicialmente?.\begin{solution} $ \left\{\begin{matrix}\tan{\left(30 \right)} = \dfrac{100}{x} \\ \tan{\left(60 \right)} = \dfrac{100}{y} \\ \end{matrix}\right. \to \left\{ x : 100 \sqrt{3}, \  y : \dfrac{100 \sqrt{3}}{3}\right\}\to 100 \sqrt{3} \land \dfrac{100 \sqrt{3}}{3} \to 115.47 \land 173.21\ m $\end{solution} \end{parts} 
    
    
\question Resuelve las siguientes ecuaciones\begin{parts} \part[1] $\cos{x}=\dfrac{1}{2}$\begin{solution} $x=60^{\circ}, x=300^{\circ}$\end{solution} \part[1] $\cos{x}=-\dfrac{1}{2}$\begin{solution} $x=120^{\circ}, x=240^{\circ}$\end{solution} \part[1] $4(\cos{x})^2-3=0$\begin{solution} $x=30^{\circ}, x=150^{\circ}, x=210^{\circ}, x=330^{\circ}$\end{solution} \part[1] $4(\cos{x})^2-2=0$\begin{solution} $x=45^{\circ}, x=135^{\circ}, x=225^{\circ}, x=315^{\circ}$\end{solution} \end{parts}

\addpoints

\end{questions}

\end{document}
%\grid
